\section{Derived Metric Analysis Pipeline}
\label{sec:derived_metric_analysis}

The analysis pipeline operates on derived metric tables generated by an upstream activation-analysis stage. 
The upstream stage computes primitive hidden-state and token-distribution metrics from transformer activations, while the present stage performs secondary evidence aggregation, condition-effect estimation, alpha-response regression, and layerwise transition analysis.

\subsection{Primitive Metrics Computed Upstream}
\label{subsec:primitive_metrics}

Let $h_{\ell}^{(t)} \in \mathbb{R}^{d}$ denote the residual-stream hidden state at layer $\ell$ and token position $t$, and let $v_{\mathrm{att}} \in \mathbb{R}^{d}$ denote the latent attractor direction. 
The upstream pipeline defines directional geometrical capture as:

\begin{equation}
\mathrm{Cos}(\theta_{\ell}^{(t)})
=
\frac{
h_{\ell}^{(t)} \cdot v_{\mathrm{att}}
}{
\|h_{\ell}^{(t)}\|_2 \, \|v_{\mathrm{att}}\|_2
}.
\end{equation}

The distance to the reference endpoint $h_{\mathrm{ref}}$ is computed as:

\begin{equation}
\mathcal{D}_{L2}
\left(
h_{\ell}^{(t)}, h_{\mathrm{ref}}
\right)
=
\left\|
h_{\ell}^{(t)} - h_{\mathrm{ref}}
\right\|_2.
\end{equation}

For token-distribution uncertainty, let $P(x_i)$ denote the softmax probability of vocabulary item $x_i \in V$. 
The mean token entropy is:

\begin{equation}
H(X)
=
-
\sum_{i \in V}
P(x_i)
\log_2
P(x_i).
\end{equation}

For alpha-steering interventions, the idealized logit-margin response is modeled as:

\begin{equation}
\Delta \mathcal{L}
=
\beta \alpha + \mathcal{C},
\end{equation}

where $\alpha$ denotes the intervention strength, $\beta$ the causal propagation slope, and $\mathcal{C}$ the intercept term.

The present analysis script does not recompute these primitive quantities from raw activations. 
Instead, it consumes derived metric columns such as cosine alignment, attractor projection, $L_2$ distance, entropy, log-probability, intervention delta, and alpha-response values.

\subsection{Derived Metric Tables}
\label{subsec:derived_tables}

Let $\mathcal{F} = \{f_1,\ldots,f_N\}$ denote the collection of derived metric tables. 
Each row $r$ contains grouping variables $g(r)$ and numerical metric values $m_j(r)$:

\begin{equation}
r
=
\left(
g(r),
m_1(r),
m_2(r),
\ldots,
m_K(r)
\right).
\end{equation}

The grouping key may include experimental condition, reference condition, base condition, intervention name, layer, layer band, module, intervention sign, token step, or alpha value:

\begin{equation}
g(r)
\subseteq
\{
\mathrm{condition},
\mathrm{reference\_condition},
\mathrm{base\_condition},
\mathrm{intervention},
\ell,
\mathrm{layer\_band},
\mathrm{module},
\mathrm{sign},
t,
\alpha
\}.
\end{equation}

The metric set may include:

\begin{equation}
m_j
\in
\{
\mathrm{projection},
\mathrm{cosine},
\mathcal{D}_{L2},
H,
\log P,
\Delta,
|\Delta|,
\mathrm{behavioral\ audit\ metrics}
\}.
\end{equation}

Only finite numerical values are included in statistical summaries.

\subsection{Global Metric Aggregation}
\label{subsec:global_aggregation}

For each metric $m$, the finite-sample count is:

\begin{equation}
n_m
=
\sum_{r}
\mathbf{1}
\left[
\mathrm{finite}
\left(
m(r)
\right)
\right].
\end{equation}

The global mean is:

\begin{equation}
\mu_m
=
\frac{1}{n_m}
\sum_{r:\,\mathrm{finite}(m(r))}
m(r).
\end{equation}

The sample standard deviation is:

\begin{equation}
s_m
=
\sqrt{
\frac{1}{n_m - 1}
\sum_{r:\,\mathrm{finite}(m(r))}
\left(
m(r) - \mu_m
\right)^2
}.
\end{equation}

The minimum and maximum are:

\begin{equation}
m_{\min}
=
\min_{r:\,\mathrm{finite}(m(r))}
m(r),
\qquad
m_{\max}
=
\max_{r:\,\mathrm{finite}(m(r))}
m(r).
\end{equation}

The approximate $95\%$ confidence interval for the mean is reported as:

\begin{equation}
\mathrm{CI}_{95}(m)
=
\left[
\mu_m - 1.96 \frac{s_m}{\sqrt{n_m}},
\;
\mu_m + 1.96 \frac{s_m}{\sqrt{n_m}}
\right].
\end{equation}

\subsection{Grouped Metric Aggregation}
\label{subsec:grouped_aggregation}

For a grouping key $g$, the grouped finite count is:

\begin{equation}
n_{g,m}
=
\sum_{r \in g}
\mathbf{1}
\left[
\mathrm{finite}
\left(
m(r)
\right)
\right].
\end{equation}

The grouped mean is:

\begin{equation}
\mu_{g,m}
=
\frac{1}{n_{g,m}}
\sum_{r \in g:\,\mathrm{finite}(m(r))}
m(r).
\end{equation}

The grouped sample standard deviation is:

\begin{equation}
s_{g,m}
=
\sqrt{
\frac{1}{n_{g,m} - 1}
\sum_{r \in g:\,\mathrm{finite}(m(r))}
\left(
m(r)-\mu_{g,m}
\right)^2
}.
\end{equation}

The grouped extrema are:

\begin{equation}
m^{\min}_{g}
=
\min_{r \in g:\,\mathrm{finite}(m(r))}
m(r),
\qquad
m^{\max}_{g}
=
\max_{r \in g:\,\mathrm{finite}(m(r))}
m(r).
\end{equation}

\subsection{Condition-Level Effects}
\label{subsec:condition_effects}

For each metric $m$, the condition-level effect is defined as the difference between a target or intervention condition $c_t$ and a baseline condition $c_b$:

\begin{equation}
\Delta_m(c_t,c_b)
=
\mu_{c_t,m}
-
\mu_{c_b,m}.
\end{equation}

The relative condition effect is:

\begin{equation}
\Delta^{\mathrm{rel}}_m(c_t,c_b)
=
\frac{
\Delta_m(c_t,c_b)
}{
|\mu_{c_b,m}|
},
\qquad
\mu_{c_b,m} \neq 0.
\end{equation}

For directional cosine and attractor projection metrics, a positive condition effect indicates stronger alignment with the latent attractor direction:

\begin{equation}
\Delta_m(c_t,c_b) > 0,
\qquad
m \in
\{
\mathrm{cosine},
\mathrm{projection}
\}.
\end{equation}

For $L_2$ distance metrics, a negative condition effect indicates compression toward the reference endpoint:

\begin{equation}
\Delta_m(c_t,c_b) < 0,
\qquad
m \in
\{
\mathcal{D}_{L2}
\}.
\end{equation}

For entropy metrics, a positive condition effect indicates entropy amplification:

\begin{equation}
\Delta_m(c_t,c_b) > 0,
\qquad
m \in
\{
H
\}.
\end{equation}

\subsection{Alpha-Steering Regression}
\label{subsec:alpha_regression}

For alpha-intervention experiments, the pipeline estimates the empirical response of each derived metric $m$ to intervention strength $\alpha$ using ordinary least squares:

\begin{equation}
m(r)
=
\beta_m \alpha(r)
+
\mathcal{C}_m
+
\varepsilon_r.
\end{equation}

Let $x_r = \alpha(r)$ and $y_r = m(r)$. 
For a group of $n$ finite observations, the estimated slope is:

\begin{equation}
\hat{\beta}_m
=
\frac{
n \sum_r x_r y_r
-
\left(
\sum_r x_r
\right)
\left(
\sum_r y_r
\right)
}{
n \sum_r x_r^2
-
\left(
\sum_r x_r
\right)^2
}.
\end{equation}

The intercept is:

\begin{equation}
\hat{\mathcal{C}}_m
=
\frac{
\sum_r y_r
-
\hat{\beta}_m
\sum_r x_r
}{
n
}.
\end{equation}

The coefficient of determination is computed as:

\begin{equation}
R_m^2
=
\frac{
\hat{\beta}_m
\left(
\sum_r x_r y_r
-
\frac{
\sum_r x_r \sum_r y_r
}{n}
\right)
}{
\sum_r y_r^2
-
\frac{
\left(
\sum_r y_r
\right)^2
}{n}
}.
\end{equation}

When absolute intervention strength is available, the same regression is also estimated with:

\begin{equation}
x_r = |\alpha(r)|.
\end{equation}

The resulting slope is reported as a beta-like alpha-response coefficient:

\begin{equation}
\hat{\beta}^{|\alpha|}_m
=
\frac{
n \sum_r |\alpha(r)| y_r
-
\left(
\sum_r |\alpha(r)|
\right)
\left(
\sum_r y_r
\right)
}{
n \sum_r |\alpha(r)|^2
-
\left(
\sum_r |\alpha(r)|
\right)^2
}.
\end{equation}

\subsection{Layerwise Phase-Transition Proxy}
\label{subsec:layerwise_transition}

For each metric $m$ and layer $\ell$, the layerwise mean is:

\begin{equation}
\mu_{\ell,m}
=
\frac{1}{n_{\ell,m}}
\sum_{r:\,\ell(r)=\ell}
m(r).
\end{equation}

The adjacent layer jump is:

\begin{equation}
J_{\ell,m}
=
\mu_{\ell+1,m}
-
\mu_{\ell,m}.
\end{equation}

The maximal adjacent transition is identified by:

\begin{equation}
\ell_m^{*}
=
\arg\max_{\ell}
\left|
J_{\ell,m}
\right|.
\end{equation}

The corresponding transition magnitude is:

\begin{equation}
J_m^{*}
=
J_{\ell_m^{*},m}.
\end{equation}

The layer of maximal metric expression is:

\begin{equation}
\ell^{\mathrm{peak}}_m
=
\arg\max_{\ell}
\mu_{\ell,m},
\end{equation}

and the layer of minimal metric expression is:

\begin{equation}
\ell^{\mathrm{trough}}_m
=
\arg\min_{\ell}
\mu_{\ell,m}.
\end{equation}

The peak-to-trough range is:

\begin{equation}
R_m^{\mathrm{layer}}
=
\mu_{\ell^{\mathrm{peak}}_m,m}
-
\mu_{\ell^{\mathrm{trough}}_m,m}.
\end{equation}

The global layerwise trend is estimated by regressing layerwise means against layer index:

\begin{equation}
\mu_{\ell,m}
=
\gamma_m \ell
+
b_m
+
\varepsilon_{\ell}.
\end{equation}

The estimated layer slope is:

\begin{equation}
\hat{\gamma}_m
=
\frac{
\sum_{\ell}
\left(
\ell-\bar{\ell}
\right)
\left(
\mu_{\ell,m}-\bar{\mu}_{m}
\right)
}{
\sum_{\ell}
\left(
\ell-\bar{\ell}
\right)^2
}.
\end{equation}

The phase-transition proxy is therefore summarized by:

\begin{equation}
\mathcal{T}_m
=
\left(
\ell_m^{*},
J_m^{*},
\ell^{\mathrm{peak}}_m,
\ell^{\mathrm{trough}}_m,
R_m^{\mathrm{layer}},
\hat{\gamma}_m
\right).
\end{equation}

\subsection{Evidence Table Construction}
\label{subsec:evidence_table}

The final evidence table combines four classes of secondary evidence:

\begin{equation}
\mathcal{E}
=
\mathcal{E}_{\mathrm{global}}
\cup
\mathcal{E}_{\mathrm{condition}}
\cup
\mathcal{E}_{\alpha}
\cup
\mathcal{E}_{\mathrm{layer}}.
\end{equation}

Here,

\begin{equation}
\mathcal{E}_{\mathrm{global}}
=
\left\{
\mu_m,
s_m,
m_{\min},
m_{\max},
\mathrm{CI}_{95}(m)
\right\},
\end{equation}

\begin{equation}
\mathcal{E}_{\mathrm{condition}}
=
\left\{
\Delta_m(c_t,c_b),
\Delta^{\mathrm{rel}}_m(c_t,c_b)
\right\},
\end{equation}

\begin{equation}
\mathcal{E}_{\alpha}
=
\left\{
\hat{\beta}_m,
\hat{\mathcal{C}}_m,
R_m^2
\right\},
\end{equation}

and

\begin{equation}
\mathcal{E}_{\mathrm{layer}}
=
\left\{
\ell_m^{*},
J_m^{*},
\ell^{\mathrm{peak}}_m,
\ell^{\mathrm{trough}}_m,
R_m^{\mathrm{layer}},
\hat{\gamma}_m
\right\}.
\end{equation}

This secondary analysis does not by itself establish a full mechanistic causal claim. 
It provides evidence summaries over derived metrics. 
A stronger Latent Discourse Regime claim requires that geometric alignment, specificity against controls, alpha-response structure, behavioral coupling, and replication evidence jointly pass their respective validation gates.